% Idioma y codificación
% \usepackage[utf8]{inputenc}
\usepackage[T1]{fontenc}
\usepackage[spanish,es-tabla]{babel}

% Tipografía y estilos de texto
\usepackage{times}          % Fuente Times New Roman
\usepackage{setspace}       % Ajuste de espaciado entre líneas
\usepackage{textcomp}       % Símbolos adicionales
\usepackage{gensymb}        % Símbolos genéricos (por ejemplo, ° para grados)
\usepackage{enumitem}       % Personalización de listas enumeradas
\usepackage{multicol}
% \usepackage{lipsum} 
\usepackage{comment}        % Para el entorno comment

% Matemáticas
\usepackage{amssymb, amsmath}   % Símbolos y entornos matemáticos
\usepackage{bm}                 % Negritas en matemáticas
% \usepackage{empheq}             % Marcos y resaltado de ecuaciones
% \usepackage{cancel}         % Cancelación de términos en ecuaciones
% \usepackage{siunitx}        % Unidades del SI

% Gráficos y colores
\usepackage{graphicx}               % Inclusión de gráficos
\usepackage{subcaption}             % Subfiguras
\usepackage[table,xcdraw]{xcolor}   % Opciones avanzadas de color
% \usepackage[export]{adjustbox}      % Ajustar alineación de gráficos
% \usepackage{wrapfig}              % Ajustar texto alrededor de imágenes
% \usepackage{pgfplots}             % Graficar funciones directamente 
\usepackage{float}          % Control preciso de la posición de tablas/figuras
% \usepackage{tcolorbox}      % Cajas de texto con formato

% Formato y estructura del documento
% \usepackage{vmargin}      % Márgenes personalizados
\usepackage{geometry}       % Márgenes personalizados
\usepackage{fancyhdr}       % Encabezados y pies de página personalizados
\usepackage{tocloft}        % Personalización de la tabla de contenidos
\usepackage{pdflscape}      % Para rotar la página
\usepackage{rotating}       % Permite rotar figuras y tablas

% Hipervínculos y navegación
\usepackage[hidelinks]{hyperref}    % Enlaces sin color ni marco visible
\usepackage{bookmark}               % Mejora la navegación en el PDF
\usepackage{titlesec}       % Para redefinir títulos y secciones
\usepackage{etoolbox}
\usepackage{multibib}
\newcites{A}{Bibliografía}
\newcites{B}{Referencias}

% Tablas avanzadas
\usepackage{array}          % Mejoras en tablas
\usepackage{multirow}       % Combinar filas en tablas
\usepackage{tabularray}     % Alternativa moderna a tabular
\UseTblrLibrary{booktabs}
\SetTblrDefault{conthead = {}}


